% ============================================================
%  chapters/app-b-projection-theory.tex
%  Appendix B: Projection Theory, Information Loss, Surjective Geometry
% ============================================================

\section*{B.1 Introduction}

The previous appendix established that symbolic systems emerge through
equivalence relations and quotient constructions. We now study the geometry
of projection itself.

\begin{tcolorbox}[enhanced, colback=RSVPdeep, colframe=RSVPdeep,
  coltext=white, arc=6pt, boxrule=0pt, left=6mm, right=6mm, top=5mm, bottom=5mm]
  \textbf{Central claim.} Every useful representation is necessarily
  many-to-one. A representation that loses nothing is not a representation
  at all.
\end{tcolorbox}

% ── B.2–B.3 Projections and Fibers ────────────────────────────
\section*{B.2 Projections}

\begin{definition}{Projection}{projection-app}
  Let $X$ and $M$ be sets. A \emph{projection} is a function
  $\pi : X \to M$. The space $X$ is the \emph{source space};
  $M$ is the \emph{representation space}.
\end{definition}

\begin{example}
  A camera image projects $\mathbb{R}^3 \to \mathbb{R}^2$.
  Typing projects $G \to \Sigma$, where $G$ is finger-trajectory space
  and $\Sigma$ is the alphabet. Logical evaluation projects
  $\text{Formulae} \to \{T,F\}$.
\end{example}

\begin{definition}{Fiber}{fiber-app}
  Given $m \in M$, the \emph{fiber over $m$} is
  $F_m = \pi^{-1}(m) = \{x \in X : \pi(x) = m\}$.
\end{definition}

\begin{example}
  For $\pi(x,y) = x$, the fiber over $m = 3$ is the entire vertical line
  $F_3 = \{(3,y) : y \in \mathbb{R}\}$. The representation remembers only
  $x = 3$; everything else is forgotten.
\end{example}

% ── B.4 Information Loss Theorem ──────────────────────────────
\section*{B.4 The Information Loss Theorem}

\begin{theorem}{Information Loss}{info-loss}
  If $|X| > |M|$, then every surjective projection $\pi : X \to M$
  loses information.
\end{theorem}

\begin{proof}
  Suppose not. Then every point in $M$ would correspond to exactly one point
  in $X$, making $\pi$ injective. An injective surjection satisfies
  $|X| = |M|$, contradicting $|X| > |M|$.
\end{proof}

\begin{warning}
  Compression is not an implementation detail. It is mathematically
  unavoidable. Any projection from a richer space to a poorer one
  necessarily discards information.
\end{warning}

% ── B.5–B.6 Entropy and Chains ────────────────────────────────
\section*{B.5 Representation Entropy}

\begin{definition}{Representation Entropy}{rep-entropy}
  The \emph{representation entropy} of $m \in M$ is
  $S(m) = \log |F_m|$.
  Large fibers correspond to large uncertainty;
  small fibers to precise representations.
\end{definition}

\begin{example}
  Suppose 1000 different gestures produce the letter ``A''. Then
  $S(\texttt{A}) = \log(1000)$. The representation ``A'' contains
  less information than the original motion.
\end{example}

\section*{B.6 Projection Chains}

\begin{definition}{Projection Chain}{proj-chain}
  A \emph{projection chain} is
  $X \to M_1 \to M_2 \to \cdots \to M_n$.
\end{definition}

A canonical example:
\[
  \text{trajectory} \;\to\; \text{gesture} \;\to\;
  \text{letter} \;\to\; \text{word} \;\to\;
  \text{sentence} \;\to\; \text{semantic category.}
\]

\begin{theorem}{Entropy Monotonicity}{entropy-mono}
  Entropy cannot decrease along a projection chain.
\end{theorem}

\begin{proof}
  Fibers can only grow under composition of projections. If
  $\pi_2 \circ \pi_1$ is a composition, each final state inherits all
  ambiguity from previous stages plus additional ambiguity introduced later.
  Therefore $S_{n+1} \geq S_n$.
\end{proof}

% ── B.7–B.9 Surjective Hypercubes ────────────────────────────
\section*{B.7 The Surjective Hypercube}

\begin{definition}{Surjective Hypercube}{surj-hypercube}
  Let $Q_n = \{0,1\}^n$. A \emph{surjective hypercube} is a pair
  $(Q_n, \pi)$ where $\pi : Q_n \to Y$ is surjective.
\end{definition}

\begin{example}
  Let $\pi(x_1,x_2,x_3,x_4) = (x_1,x_2)$. Then 16 states collapse to
  4 states, each with 4 preimages: $Q_4 \to Q_2$.
\end{example}

\section*{B.8 Hypercube Splitting}

\begin{definition}{Splitting Function}{splitting}
  A \emph{splitting function} is $\sigma : Q_n \to \{0,1\}$, partitioning
  the hypercube into $Q_n^+ = \sigma^{-1}(1)$ and $Q_n^- = \sigma^{-1}(0)$.
\end{definition}

Every Boolean decision creates a hypercube split. The penteract (Appendix~E)
is a fibered pair $(Q_n^-, Q_n^+, \pi)$ where the lower cube represents
detailed states, the upper cube compressed states, and the connecting fibers
encode admissible lifts.

% ── B.12 Fundamental Theorem ──────────────────────────────────
\section*{B.12 The Fundamental Projection Principle}

\begin{theorem}{Reality Partitioned by Representation}{reality-partition}
  Every representation partitions reality into fibers.
\end{theorem}

\begin{proof}
  Let $\pi : X \to M$. For every $m \in M$ define $F_m = \pi^{-1}(m)$.
  The collection $\{F_m\}$ forms a partition of $X$.
\end{proof}

\begin{namedthm}{The Fiber Principle}
  A letter is a fiber. A word is a fiber. A gesture is a fiber.
  A concept is a fiber. A memory is a fiber. A scientific theory is a fiber.
  Every representation is a partition imposed upon a richer reality.
\end{namedthm}

\medskip
\textit{Appendix~C} will use this machinery to derive Spherepop evaluation as
a constraint-reduction process on nested expression manifolds.
